\documentclass[11pt]{article}
\usepackage[a4paper,margin=25mm]{geometry}
\usepackage{amsmath,amssymb,amsthm,mathtools,booktabs,hyperref,xurl}
\setlength{\emergencystretch}{3em}
\newtheorem{theorem}{Theorem}
\newtheorem{lemma}[theorem]{Lemma}
\newtheorem{corollary}[theorem]{Corollary}
\newcommand{\Q}{\mathbb Q}\newcommand{\Z}{\mathbb Z}
\newcommand{\A}{\mathcal A}\newcommand{\Mat}{\operatorname{Mat}}
\title{Paired spectral differences and lower-degree\\
guarded finite-algebra encoders}
\author{\(\mathbb Z\mathbb Q\) Diophantine Project}
\date{23 November 2026}
\begin{document}\maketitle

\begin{abstract}
We improve the complete one-coordinate encoding of a guarded rational
character of a finite commutative algebra. We remove the mandatory
diagonal zero factors and pair opposite spectral differences, then use
separate bounds for the eigenvalues and the characteristic coefficients.
For algebra rank \(n\), fixed generator count, and multiplication-entry
degree \(d\), the degree bound falls from order \(n^7d\) to order
\(n^5d\). The proof includes nilpotents, repeated roots, an empty permitted
locus, and denominators vanishing at excluded characters.

Applied to the existing rational-undecidability transfers, this lowers
the degree ceilings in 82, 100, and 52 variables from
\(2^{650},2^{714},2^{500}\) to \(2^{532},2^{578},2^{425}\).
The 83-, 101-, and 53-variable tradeoffs instead have ceilings
\(2^{302},2^{349},2^{212}\). All inequalities are strict. The variable
counts and every retained arithmetic premise are unchanged; the
52- and 53-variable routes retain ordinary rational \(abc\).
\end{abstract}

\section{The algebraic theorem}

Let \(R=\Z[T_1,\ldots,T_b]\), and let \(\A\) be a commutative unital
\(R\)-algebra free of rank \(n\ge1\), with a specified basis. Choose
\(\theta_1,\ldots,\theta_m\) that generate
\(\A\otimes\Q\) over \(\Q[T_1,\ldots,T_b]\), and a guard
\(g\in\A\). The guard is required to be nonzero only at the selected
character. Put
\[
 \gamma_0=g,\qquad \gamma_i=g\theta_i\ (1\le i\le m),
 \qquad A_i=M_{\gamma_i},\qquad \ell=m+1.
\]
The notation \(M_a\) means multiplication by \(a\) in the fixed basis.
Assume all entries of the \(A_i\) have total degree at most \(d\ge0\).

\begin{theorem}\label{main}
The coefficient recipe below constructs \(F\in R[U]\) such that,
for every \(t\in\Q^b\),
\[
 \exists u\in\Q\ F(t,u)=0
 \quad\Longleftrightarrow\quad
 \exists\phi:\A_t\longrightarrow\Q\quad\phi(g)\ne0,
 \qquad \A_t=\A\otimes_{R,t}\Q.
\]
No reducedness or separability assumption is needed. With
\[
                  J=\frac{\ell n(n+1)}2,
\]
the total degree satisfies
\begin{equation}\label{degree}
 \boxed{\deg F\le
  2n(J+1)\bigl[1+d\{1+(4+6\ell)J\}\bigr].}
\end{equation}
There is exactly one extra rational coordinate, \(U\).
\end{theorem}

The earlier encoder \cite{old} used \(J_0=\ell(n+1)^2\) and the bound
\[
 2n(J_0+1)\bigl[1+d\{1+2J_0(1+\ell(J_0+1))\}\bigr].
\]
The new formula has the same logical input and output. It may replace
that encoder without changing any source equations or witnesses.

\section{Removing diagonal zeros and pairing opposite differences}

Put \(N=n+1\) and augment every multiplication matrix by one zero block:
\[
 \widetilde A_i=\operatorname{diag}(A_i,0),\qquad
 D_i=\widetilde A_i\otimes I_N-I_N\otimes\widetilde A_i.
\]

\begin{lemma}\label{paired}
There is a unique monic \(q_i(V)\in R[V]\), of degree
\(N(N-1)/2\), such that the polynomial identity
\[
                 \det(ZI_{N^2}-D_i)=Z^Nq_i(Z^2)
\]
holds in \(R[Z]\). If each entry of \(A_i\) has degree at most \(d\),
the coefficient of \(V^{N(N-1)/2-j}\) in \(q_i\) has degree at
most \(2jd\).
\end{lemma}
\begin{proof}
Work first over an algebraic closure of the fraction field of \(R\).
Triangularize \(\widetilde A_i\), with diagonal entries
\(\lambda_1,\ldots,\lambda_N\), repeated according to multiplicity.
The tensor difference has diagonal entries
\(\lambda_a-\lambda_b\). The \(N\) diagonal pairs \(a=b\) contribute
\(Z^N\). Each pair \(a<b\) contributes
\[
 (Z-\lambda_a+\lambda_b)(Z+\lambda_a-\lambda_b)
       =Z^2-(\lambda_a-\lambda_b)^2.
\]
Thus division by \(Z^N\) has only even exponents. The original
determinant has coefficients in \(R\); deleting the identically zero
coefficients and halving the exponents therefore constructs \(q_i\)
in \(R[V]\), without any division of coefficients. The identity in
the fraction field is consequently an identity in \(R[Z]\).

The coefficient in question is the characteristic coefficient of
index \(2j\) of the matrix \(D_i\). It is a sum of principal minors
of size \(2j\), so its degree is at most \(2jd\).
\end{proof}

This argument allows further zero differences and nilpotent matrices.
It removes only the universally present \(N\) diagonal factors.
After any rational specialization, the eigenvalues of multiplication
are the geometric character values, with positive multiplicities.
For example, on a local Artinian component, multiplication by an
element is its residue value plus a nilpotent operator. Hence all
character differences, and the differences of a character value from
the extra zero, occur among the spectral differences above.

Define
\begin{equation}\label{psi}
 \psi(V)=\prod_{i=0}^m q_i(V)
        =V^J+a_1V^{J-1}+\cdots+a_J,\qquad a_0=1.
\end{equation}
Every coefficient satisfies \(\deg a_j\le2jd\). The roots of
\(\psi\) are the \emph{squares} of the relevant spectral differences.

\section{The complete integer-polynomial recipe}

Use the following spectral envelope, which deliberately contains
no characteristic coefficients:
\begin{equation}\label{envelope}
 R_0=2+N+\sum_{i=0}^m\sum_{v,w=1}^n(A_i)_{vw}^2,
       \qquad B_0=NR_0.
\end{equation}
For every \(s\in\{0,\ldots,J\}\), define
\begin{align}
 \Delta_s&=a_s^2,\qquad
 K_s=(1+\Delta_s)(2B_0)^{J-1},\label{K}\\
 C_s&=\Delta_s+(4\ell B_0+2)K_s,\label{C}\\
 M_s&=\Delta_s\sum_{i=0}^m
       C_s^{i+1}\Delta_s^{m-i}A_i,\qquad
 H_s(U)=\det(I_n-UM_s).\label{matrix}
\end{align}
Here \(J\ge1\), so all exponents are nonnegative integers.
Finally set
\begin{equation}\label{F}
 \boxed{F(T,U)=\prod_{s=0}^J
       \left[H_s(U)^2+\sum_{j=s+1}^J a_j^2\right].}
\end{equation}
Equations \eqref{psi}--\eqref{F}, the multiplication matrices, and the
identity of Lemma~\ref{paired} specify every coefficient by finite
integer-polynomial operations. An empty tail sum is zero.

\section{Separation, rationality, and all coefficient strata}

Fix a rational specialization of the base coordinates. All matrix
entries and characteristic coefficients in this section are rational
real numbers, although the eigenvalues may be complex. Every entry
has absolute value at most \(R_0\); the row-sum spectral bound shows
that every eigenvalue of \(\widetilde A_i\) has modulus at most
\(B_0\). Every difference has modulus at most \(2B_0\).

Let \(s\) be the largest index with \(a_s\ne0\), including \(s=0\)
when all the nonleading coefficients vanish. If \(s\ge1\), any
nonzero spectral difference \(\delta\) has \(\delta^2\) among
the \(s\) nonzero roots of \(\psi\). Their product has modulus
\(|a_s|\). Consequently
\[
       |\delta|\ge\frac{\sqrt{|a_s|}}{(2B_0)^{s-1}}.
\]
For every real \(t\ge0\),
\(t^2/(1+t^2)\le\sqrt t\): if \(t\le1\), use
\(t^2\le\sqrt t\), and if \(t\ge1\), use
\(t^2/(1+t^2)\le1\le\sqrt t\). Since \(2B_0>1\), this gives
\begin{equation}\label{epsilon}
  |\delta|\ge
  \frac{a_s^2}{(1+a_s^2)(2B_0)^{J-1}}
       =\frac{\Delta_s}{K_s}=:\epsilon>0.
\end{equation}
This is the step that separates the coefficient bound from the
low-degree matrix envelope.

Put \(L=C_s/\Delta_s=1+(4\ell B_0+2)/\epsilon\). Consider the vectors
\[
  v_\sigma=(\sigma(\gamma_0),\ldots,\sigma(\gamma_m))
\]
for geometric characters \(\sigma\), together with the zero vector.
For two distinct such vectors let \(h\) be the last differing index.
Their difference in coordinate \(h\) has modulus at least \(\epsilon\).
All coordinate differences have modulus at most \(2B_0\). Therefore
\[
 \left|L^{h+1}(v_\sigma-v_\tau)_h\right|
       \ge\epsilon L^{h+1},\qquad
 \left|\sum_{i<h}L^{i+1}(v_\sigma-v_\tau)_i\right|
       \le2\ell B_0L^h<\epsilon L^{h+1}.
\]
The weighted sums cannot coincide. The matrix \(M_s\) is multiplication by
\[
 \beta_s=\Delta_s^{m+2}\sum_{i=0}^m L^{i+1}\gamma_i.
\]
The nonzero scalar factor does not affect separation. If
\(\sigma(g)=0\), its entire vector and its \(\beta_s\) value are zero.
If \(\sigma(g)\ne0\), the vector is nonzero. Two such characters
with the same vector agree on every generator, since
\(\sigma(\theta_i)=\sigma(\gamma_i)/\sigma(\gamma_0)\).
Thus \(\beta_s\) separates all permitted characters and sends
precisely the excluded characters to zero.

If the true index is \(s=0\), every squared difference is zero.
The extra zero eigenvalue implies that all \(\gamma_i\) have only
zero character values. There is no permitted character, all the
\(M_s\) have only zero eigenvalues, and every relevant determinant
is one.

\begin{proof}[Proof of the equivalence in Theorem~\ref{main}]
Suppose a permitted rational character \(\phi\) exists. The true
index satisfies \(s\ge1\), and \(\phi(\beta_s)\in\Q^\times\).
Taking \(u=1/\phi(\beta_s)\) makes \(H_s(u)=0\), while all later
coefficients vanish. Hence \(F(t,u)=0\).

Conversely, if \(F(t,u)=0\), one factor in \eqref{F} vanishes.
It is a sum of rational squares, so its determinant and all tail
coefficients vanish. If its \(a_s\) were zero, then
\(\Delta_s=0\), \(M_s=0\), and \(H_s(u)=1\), impossible.
Thus this factor has the true last nonzero coefficient index.
The determinant supplies a geometric character with
\(\sigma(\beta_s)=1/u\in\Q^\times\), in particular \(u\ne0\).
It has nonzero guard. Every Galois conjugate of this character has
the same rational \(\beta_s\) value; separation makes every such
conjugate identical. All generator values are therefore rational,
so \(\sigma\) is a character into \(\Q\) with nonzero guard.
\end{proof}

The reasoning uses eigenvalues with their algebraic multiplicities,
not simple roots or one-dimensional eigenspaces. Repeated roots and
nilpotents are therefore covered. If a source denominator vanishes
at an excluded character, that character still maps to zero under
the separator. There is no need to invert its denominator.

\section{The degree calculation}

Lemma~\ref{paired} and the product \eqref{psi} give
\(\deg a_j\le2jd\). From \eqref{envelope}--\eqref{C},
\[
 \deg R_0,\deg B_0\le2d,\qquad
 \deg\Delta_s\le4Jd,\qquad
 \deg K_s\le(6J-2)d,\qquad
 \deg C_s\le6Jd.
\]
The largest entry-degree bound in \eqref{matrix} is obtained at
\(i=m\), and is
\[
                    d_\beta=d+(4+6\ell)Jd.
\]
Consequently \(\deg H_s\le n(1+d_\beta)\). The square of each
tail coefficient has degree at most \(4Jd\), which is bounded by
\(2n(1+d_\beta)\). There are \(J+1\) factors in \eqref{F};
this proves \eqref{degree}, including \(d=0\).

\section{Application to the current undecidability transfers}

The sources, rational coordinates, and finite algebras are exactly
those of the common-numerator and finite-wrapper papers
\cite{common,wrapper}. In particular every one of the nine integer
predicates remains present. The final polynomial is obtained by
applying \eqref{F} to their already specified multiplication algebra.
No coefficient in a source DAG is changed. The integer instance is
encoded in coefficients and is not a quantified rational coordinate.

The finite-wrapper algebras have ranks \(2^{56},2^{65},2^{35}\)
and \(\ell=54,63,33\), respectively. Their multiplication-entry
degree bounds are those independently replayed in the source ledger.
For clarity, let
\[
 \delta=15616+233856\cdot4+
           (233952\cdot4+467712\cdot4^2)9^{59}.
\]
This is the retained universal-pair \((58,4)\) source bound. If
\(d_0=d_g+d_\theta\) is the common-numerator operator bound, the
finite-wrapper bound used here is
\[
                         d=d_0(2\delta+1)+\delta.
\]
The three values of \(d_0\) are
\(2302058,4126214,3309916\). Substitution in \eqref{degree}
gives the first three rows below.

\begin{center}\small
\begin{tabular}{lrrr}
\toprule
Construction & Rational variables & Old ceiling & New ceiling\\
\midrule
Finite wrapper, general &82&\(2^{650}\)&\(2^{532}\)\\
Finite wrapper, specified &100&\(2^{714}\)&\(2^{578}\)\\
Finite wrapper, ordinary \(abc\)&52&\(2^{500}\)&\(2^{425}\)\\
Common numerator, general &83&\(2^{413}\)&\(2^{302}\)\\
Common numerator, specified &101&\(2^{477}\)&\(2^{349}\)\\
Common numerator, ordinary \(abc\)&53&\(2^{263}\)&\(2^{212}\)\\
Earlier shared points, general &99&\(2^{526}\)&\(2^{384}\)\\
Earlier shared points, specified &117&\(2^{591}\)&\(2^{430}\)\\
Earlier shared points, ordinary \(abc\)&69&\(2^{377}\)&\(2^{278}\)\\
\bottomrule
\end{tabular}
\end{center}
All displayed ceilings are strict upper bounds, not exact polynomial
degrees. For the common-numerator and earlier shared-point routes,
the final degree bound is
\[
                     \max\{4\mathcal D,2\delta\},
\]
where \(\mathcal D\) is \eqref{degree}. The factor four accounts
for the ternary nonnegative substitution and the final sum of squares.
The \(53\)-variable encoder contribution alone is below \(2^{196}\),
but its source contribution forces the stated ceiling \(2^{212}\).
Omitting this source floor would give an invalid degree claim.

All general and specified rows retain the arithmetic premises of
their integer predicates. The general and ordinary-\(abc\) sufficient
fixed parameters remain ineffective. The specified route retains
its fixed \(L=143\) member. The new algebraic theorem introduces no
additional arithmetic hypothesis. These transfers do not depend on
the separate unresolved emended Jones source \(U_{\rm em}\).

The current single-input constructions in 25, 27, and 11 variables
already use stronger, specialized point and terminal-root encoders.
Their v88 degree bounds remain
\[
 48{,}318{,}382{,}080,\quad183{,}368{,}679{,}424,\quad
 3{,}261{,}569,
\]
with \(16{,}460{,}929\) and \(9{,}808{,}001\) for the two weak-\(abc\)
11-variable variants. The generic theorem here does not improve
those five specialized bounds, and no change to them is claimed.

\section{Exact specification and verification scope}

The directory \nolinkurl{undecidability/alternative_v93} contains
\nolinkurl{paired_spectral.py}, implementing every step of
\eqref{psi}--\eqref{F}, and
\nolinkurl{check_paired_spectral.py}, an independent spectrum and
degree-ledger checker. Its exact report is
\nolinkurl{certificate.json}. The sources and their certificates
are pinned by SHA-256. The checker replays multiplication-entry
degree bounds from the full common-numerator and shared source DAGs
and checks the finite-wrapper degree rescaling independently.

Small exact examples compare the computed squared-difference
polynomial with the independent unordered-pair product of known
spectra. They include nonsplit quadratic fields, nonreal spectra,
zero terminal roots, independent nilpotents in a nonmonogenic
algebra, a nilpotent guard excluding every character, and a denominator
vanishing on excluded nonreduced components. A symbolic family
checks the diagonal quotient and even-power identity before
specialization. Every wrong coefficient stratum is checked exactly.
These finite tests support the coefficient construction; the
universal equivalence is the proof above. The enormous final
polynomials and multiplication matrices have not been expanded.

\begin{thebibliography}{9}
\bibitem{old}
\(\Z\Q\) Diophantine Project,
\emph{One rational coordinate for a finite family of algebraic witnesses},
23 November 2026.
\nolinkurl{undecidability/papers/joint_finite_algebra_encoder.tex}.
\bibitem{common}
\(\Z\Q\) Diophantine Project,
\emph{A common numerator point for simultaneous integer certificates
and rational undecidability in 83, 101, and 53 variables},
23 November 2026.
\nolinkurl{undecidability/papers/common_numerator_transfers.tex}.
\bibitem{wrapper}
\(\Z\Q\) Diophantine Project,
\emph{Absorbing the mixed-integer wrapper into a finite algebra:
Rational undecidability in 82, 100, and 52 variables},
23 November 2026.
\nolinkurl{undecidability/papers/finite_wrapper_transfer.tex}.
\bibitem{v88}
\(\Z\Q\) Diophantine Project,
\emph{Product poles and lower degrees in the few-variable constructions},
23 November 2026.
\nolinkurl{papers/few_variables25/bb_product_degree_transfer_20261123_v88.tex}.
\end{thebibliography}
\end{document}
